\documentclass[11pt,a4paper]{article}
\usepackage[T1]{fontenc}
\usepackage{lmodern,amsmath,amssymb,amsthm,booktabs,longtable,array,microtype}
\usepackage[margin=27mm]{geometry}
\usepackage[hidelinks]{hyperref}
\newcommand{\F}{\mathbb F_2}
\newcommand{\wt}{\operatorname{wt}}
\newcommand{\rank}{\operatorname{rank}}
\newtheorem{theorem}{Theorem}
\newtheorem{lemma}[theorem]{Lemma}
\newtheorem{proposition}[theorem]{Proposition}
\title{A reproducible spectral certificate for the nonexistence of\newline homogeneous quartic rotation symmetric bent functions\newline in twelve variables}
\author{Filipe Martins Vono}
\date{16 October 2026}
\begin{document}
\maketitle
\begin{abstract}
We give a computer-assisted certificate excluding bent functions from the
43-dimensional space of homogeneous quartic rotation symmetric Boolean
functions in twelve variables. The argument uses the restrictions
$g_z(u)=f(u,u+z)$, necessary energy identities at all 64 fiber frequencies,
and finite Poisson constraints. A triangular parametrization covers all
$2^{43}$ coefficient vectors without evaluating each complete truth table.
The final two stages retain 128 classes and then none. Two implementations,
one using signed spectral transport and one evaluating each fiber directly,
reproduce the same stage counts. We state the necessity of every family of
filters and distinguish exact computation from the auxiliary statements
checked in Lean. The result concerns twelve variables; it does not settle
the conjecture in arbitrary dimension.
\end{abstract}

\section{Scope and conventions}
The nonexistence conjecture for homogeneous rotation symmetric bent
functions of degree greater than two is discussed in
\cite{SM,MCF,Sun}. This paper addresses the quartic case in twelve
variables with a reproducible computational certificate. We do not claim
priority over all earlier searches or claim a new general obstruction
for every even dimension. The mathematical and executable descriptions
below are intended to make the finite result independently assessable.

Write $R_n$ for left rotation on $n$ bits. Integers encode vectors with
coordinate zero in the least significant bit, and addition means XOR.
For a Boolean function $f:\F^n\to\F$, put
\[
 W_f(a)=\sum_{x\in\F^n}(-1)^{f(x)+\langle a,x\rangle}.
\]
For even $n$, bentness means $|W_f(a)|=2^{n/2}$ for every $a$.
In this paper $n=12$, $u,z,b\in\F^6$, and
\[
 g_z(u)=f(u,u+z),\quad W_z(b)=W_{g_z}(b),\quad A_z=W_z(0).
\]
Let $\mathcal H$ be the binary span of the rotation orbits of squarefree
degree-four monomials. Orbits are ordered by the lexicographically smallest
support tuple. There are 43 orbits: forty of size 12, two of size 6, and
one of size 3. Thus $\dim\mathcal H=43$ and $|\mathcal H|=2^{43}$.
The zero function is included and is harmless for the exclusion statement.

\begin{theorem}\label{thm:main}
No function in $\mathcal H$ is bent.
\end{theorem}
The proof consists of the necessary conditions in Sections 2--4, exhaustive
enumeration in the triangular coordinates of Section 5, and the execution
records supplied with the source. This is a computer-assisted theorem,
not an end-to-end Lean formalization.

\section{Fiber identities and signed transport}
For any $\alpha,\beta\in\F^6$, changing variables from $(x,y)$ to
$(u,z)=(x,x+y)$ gives
\begin{equation}\label{eq:walsh}
 W_f(\alpha,\beta)=\sum_z(-1)^{\langle\beta,z\rangle}W_z(\alpha+\beta).
\end{equation}
Consequently, if $f$ is bent, Parseval in the variable $z$ gives
\begin{equation}\label{eq:energy}
 \sum_z W_z(b)^2=4096\qquad(b\in\F^6).
\end{equation}
For any set $S$ of fibers, the partial energy
$E_S(b)=\sum_{z\in S}W_z(b)^2$ is therefore at most 4096.
These 64 identities are necessary for bentness; they are not a substitute
for testing all 4096 coefficients of $W_f$.

Rotation by six positions interchanges the two blocks, so
$g_z(u+z)=g_z(u)$. Pairing $u$ with $u+z$ proves
\begin{equation}\label{eq:zero}
 \langle b,z\rangle=1\quad\Longrightarrow\quad W_z(b)=0.
\end{equation}
For $0\le k<6$, define $\tau_k(z)$ as the low six bits of $R_{12}^k(0,z)$.
Rotation sends $(u,u+z)$ to
$(R_6^ku+\tau_k(z),R_6^ku+\tau_k(z)+R_6^kz)$. Hence
\begin{equation}\label{eq:transport}
 W_{R_6^kz}(b)=(-1)^{\langle b,\tau_k(z)\rangle}
                 W_z(R_6^{-k}b).
\end{equation}
The sign cannot be omitted. Squared energies do not detect its omission,
but the Poisson sums below do. In particular, $A_z$ is constant on each
cyclic orbit of $z$.

\section{Finite algebraic foundations}
All finite checks in this section are performed independently by
\texttt{audit\_foundations.py}; their explicit orbit table and values are
in \texttt{results/foundations.json}. These checks use only the 43
generators or small restriction spaces, not a random sample of $\mathcal H$.

\begin{lemma}\label{lem:div}
For every $f\in\mathcal H$, $8\mid W_z(b)$ for all $z,b$.
Moreover, $16\mid A_z$ for $z=9,27,63$,
$W_9(63)=W_{63}(63)=0$, and $16\mid W_{27}(63)$.
\end{lemma}
\begin{proof}
The algebraic normal form of every generator restricted to every fiber
has degree at most three, as verified by the binary M\"obius transform.
Degree at most three is preserved by binary linear combinations. For
$z\ne0$, translation invariance makes $g_z$ a function on the
five-dimensional quotient by $\langle z\rangle$. If $\langle b,z\rangle=0$,
the same holds for $g_z+\langle b,u\rangle$. A polynomial of degree less
than five on $\F^5$ has even weight: summing each monomial over the cube
gives zero modulo two. Its Walsh sum is divisible by four, and lifting
to $\F^6$ doubles this sum. Otherwise use \eqref{eq:zero}.

At $z=0$, only the three short monomial orbits survive; their restrictions
are $Q_1,Q_2,Q_3$, where
\[
 Q_d=\sum_{i=0}^5u_i u_{i+d}\ (d=1,2),\qquad
 Q_3=u_0u_3+u_1u_4+u_2u_5.
\]
The eight combinations have all their Walsh coefficients divisible by eight.
This is a finite check on $8\cdot64$ coefficients.

For $z=9,27$, the induced affine map is
$T_z(u)=R_6^3u+\tau_3(z)$, with \emph{correct} constants
$\tau_3(9)=1$ and $\tau_3(27)=3$. It satisfies $T_z^2(u)=u+z$,
so its sixteen orbits all have size four. The value of $g_z$ is constant
on each orbit. Each of the 43 generators is one on an even number of
these orbits. This parity condition is linear in the generator values,
so it holds for every combination. Thus $8\mid\wt(g_z)$ and
$16\mid A_z$. For $z=9$, $T_z$ reverses the parity of $u$, giving
$W_9(63)=0$. For $z=27$, that parity is constant on each orbit and equals
one on eight of the sixteen orbits; the same even-weight argument gives
$16\mid W_{27}(63)$.

For $z=63$, the negacyclic shift $u\mapsto R_6u+1$ has five orbits
of length 12 and one of length 4, represented by $0,2,4,6,10,12$.
Let their values be $v_1,\ldots,v_6$. Direct evaluation on each of
the 43 generators gives $v_1+\cdots+v_6=0$ in $\F$; linearity
extends this to every $f$. If $k=\sum_{i=1}^5v_i$ and $\delta=v_6$
as integers, $k\equiv\delta\pmod2$ and
\[
 \wt(g_{63})=12k+4\delta=24q+16\delta,
 \qquad A_{63}\in\{-64,-32,-16,16,32,64\}.
\]
The negacyclic shift reverses parity, so $W_{63}(63)=0$.
\end{proof}
For redundancy, enumerating the complete images of the three restriction
maps gives ranks $9,10,5$ for $z=9,27,63$, respectively, and confirms
the assertions on all $512+1024+32$ image functions.

\begin{proposition}\label{prop:sector}
If $f\in\mathcal H$ is bent, $16\mid A_0$ and
$c_{40}+c_{26}c_{42}=1$ in $\F$, where indices are zero-based.
\end{proposition}
\begin{proof}
The nonzero fiber orbits have representatives
$1,3,5,7,9,11,13,15,21,23,27,31,63$. Only $9,27,63$ have odd
orbit sizes. Divide \eqref{eq:energy} at $b=0$ by 64 and reduce
modulo two; Lemma~\ref{lem:div} leaves $(A_0/8)^2=0$ modulo two.
At $b=63$, all odd-weight fibers vanish. The even-weight fibers
$9$ and $63$ also vanish. The only other nonzero orbits have
representatives $3,5,15,23,27$; their total energies are divisible by
$384$ (size six times $8^2$, or size three times $16^2$).
Thus $W_0(63)^2\equiv4096\equiv256\pmod{384}$.
The eight values for $(c_{26},c_{40},c_{42})$ in lexicographic order are
\[
\begin{array}{c|rrrrrrrr}
(c_{26},c_{40},c_{42})&000&001&010&011&100&101&110&111\\\hline
A_0&64&8&0&16&16&32&0&-8\\
W_0(63)&0&-8&16&-16&0&-32&-32&8
\end{array}
\]
and give the stated equivalence. Exactly four sectors remain.
\end{proof}

\section{Necessity of all sieve filters}
\subsection{Subspace Poisson constraints}
For any binary subspace $V\le\F^6$, any $b,\gamma\in\F^6$, define
$L_V(b,\gamma)=\sum_{z\in V}(-1)^{\langle\gamma,z\rangle}W_z(b)$.
Fourier inversion of \eqref{eq:walsh} gives
\begin{equation}\label{eq:poisson}
 L_V(b,\gamma)=\frac{|V|}{64}
   \sum_{\beta\in\gamma+V^\perp}W_f(b+\beta,\beta).
\end{equation}
For bent $f$, each summand is $64$ times a sign. Rotation symmetry
forces identical signs on each cyclic orbit of the 12-bit frequency.
If the coset has orbit multiplicities $m_1,\ldots,m_t$, necessarily
\begin{equation}\label{eq:allowed}
 L_V(b,\gamma)\in
 |V|\left\{\sum_{j=1}^t m_j\epsilon_j:\epsilon_j\in\{-1,1\}\right\}.
\end{equation}
Both programs construct this finite set by repeated signed addition.
Allowing independent signs in separate tests enlarges the feasible set;
it cannot discard a bent function. The code checks closure under XOR
for each subspace before constructing its table.

The 24 spaces used are the thirteen lines $\langle r\rangle$ for the
nonzero representatives listed above, nine further spaces in the table,
and the final two spaces $V_3'$ and $V_4$.
\begin{center}\small
\begin{tabular}{ll}
\toprule
Code name & Elements or basis (span over $\F$)\\\midrule
\texttt{V2\_21} & $\{0,21,42,63\}$\\
\texttt{V2\_27} & $\{0,27,45,54\}$\\
\texttt{V2\_9a}, \texttt{V2\_9b} & $\langle9,18\rangle$, $\langle9,54\rangle$\\
\texttt{V2\_7a}, \texttt{V2\_7b} & $\langle7,56\rangle$, $\langle7,14\rangle$\\
\texttt{V2\_7c}, \texttt{V2\_7d} & $\langle7,21\rangle$, $\langle7,28\rangle$\\
\texttt{V3} & $\langle9,18,36\rangle$\\
$V_3'$ (\texttt{V3\_7a}) & $\langle63,9,7\rangle$\\
$V_4$ (\texttt{V4}) & $\langle21,42,9,18\rangle$\\\bottomrule
\end{tabular}
\end{center}
In particular, $V_2=\langle21,42\rangle$ and $V_3=\langle9,18,36\rangle$
are not nested: $V_2\cap V_3=\langle63\rangle$ and $V_2+V_3=V_4$.

\subsection{Energy congruences}
For each cyclic set $S$ of fibers processed by the sieve, the following
congruences are used only at the indicated stages.
\begin{enumerate}
\item After the initial fibers $0,63,21,42$, the residual energy at
$b=0$ is divisible by 384. The unprocessed size-six orbits contribute
multiples of $6\cdot64$, and the size-three orbits contribute
multiples of $3\cdot256$ by Lemma~\ref{lem:div}.
\item Once $9,18,36,27,54,45$ have also been processed, every
remaining fiber orbit has length six. Residual energies at $b=0,63$
are divisible by 384, since these frequencies are fixed by rotation
and squares in \eqref{eq:transport} agree. The initial $b=63$ test
is the stronger sector congruence in Proposition~\ref{prop:sector}.
\item At $b=21$, $R_6^2b=b$. Outside $\{0,21,42,63\}$, the
$R_6^2$-orbits have size three. Therefore residual energy is divisible
by $3\cdot64=192$ at every stage after the initial one.
\item After processing the size-three fiber orbits, residual energy
summed over the three frequencies $\{9,18,36\}$ or $\{27,54,45\}$
is divisible by 192. Indeed rotation acts on the remaining pairs
$(z,b)$ in length-six orbits, and each squared term is divisible by 64;
this even gives divisibility by 384.
\end{enumerate}
The size-three contribution at $b=21$ can be computed as $3W_{27}(21)^2$.
Indeed, $R_6^{-2}21=21$, so two successive rotations preserve the
squared coefficient in \eqref{eq:transport}. On a three-element
fiber orbit, this implies equality of all three squares.
This explains the prefilter in the direct implementation.

\subsection{Autocorrelation and diagonal Walsh congruences}
Equation~\eqref{eq:walsh} at $b=0$ says that the Fourier transform of
$A$ consists of the diagonal coefficients $W_f(\beta,\beta)$.
If $f$ is bent, inversion gives
$\sum_z A_zA_{z+63}=0$. Dividing by two and grouping complementary
orbits yields the initial necessary condition
\[
 A_0A_{63}+A_{21}^2\equiv0\pmod3.
\]
Once the size-three orbits are known it yields
\begin{equation}\label{eq:ac}
 A_0A_{63}+A_{21}^2+3A_9A_{27}\equiv0\pmod{192}.
\end{equation}
For distinct complementary size-six orbits the residual contribution
is $6A_rA_s$; for a self-complementary orbit it is $3A_r^2$.
Both are divisible by 192.

The diagonal Walsh coefficients at $\beta=0,63$ are $\pm64$.
All omitted terms in their cyclic orbit sums are divisible by 48:
size-six orbits have $8\mid A_r$, and size-three orbits have
$16\mid A_r$. Hence the initial tests are
\begin{equation}\label{eq:w48}
 A_0+A_{63}\pm2A_{21}\equiv16\ \text{or}\ 32\pmod{48}.
\end{equation}
These tests use $\beta=0,63$; $b=21$ in the energy tests is a different
frequency variable.

\section{Exhaustive triangular enumeration}
Let $B$ be the $43\times4096$ matrix whose rows are generator truth
tables, with columns ordered by $(z,u)$, grouping complete fiber
orbits as in Table~\ref{tab:stages}. Row reduction over $\F$ gives
$M=PB$, with $P$ invertible and 43 pivots. The direct implementation
explicitly checks $PB=M$ and that the pivot submatrix of $M$ is the
identity. Thus it also checks injectivity of the coefficient map.

Group rows by the stage containing their pivot. A later group is zero
on all earlier fiber columns. Enumerate all binary combinations of the
first group, filter, and extend each survivor by all combinations of
the next group. Every original coefficient vector has exactly one
such sequence of extensions. No future choice can change a quantity
already tested, and every test is necessary for bentness. Therefore a
zero terminal count proves Theorem~\ref{thm:main}.

\begin{table}[htbp]\centering
\caption{Exact counts. Counts refer to classes of restrictions at the
stated rank, not equally sized lists of complete functions.}\label{tab:stages}
\begin{tabular}{lrrr}\toprule
New fiber orbit(s) & Rank & Rank increment & Survivors\\\midrule
$0,63,21,42$: initial congruences &12&12&712\\
Same fibers: energy and Poisson &12&0&361\\
$\mathcal O_9$ &18&6&20,158\\
$\mathcal O_{27}$ &24&6&212,296\\
$\mathcal O_7$ &27&3&8,688\\
$\mathcal O_1$ &31&4&26,496\\
$\mathcal O_{11}$ &34&3&94,936\\
$\mathcal O_{13}$ &36&2&149,448\\
$\mathcal O_{31}$ &37&1&106,552\\
$\mathcal O_3$ &41&4&349,680\\
$\mathcal O_5$ &41&0&23,636\\
$\mathcal O_{15}$ &42&1&128\\
$\mathcal O_{23}$ &43&1&0\\\bottomrule
\end{tabular}
\end{table}

At rank 12 the initial conditions are the energy bounds at $b=0,21$,
their residual congruences, the sector congruence at $b=63$, the
autocorrelation condition modulo three, and \eqref{eq:w48}. The next
filter adds all 64 energy bounds and Poisson on
$\langle63\rangle,\langle21\rangle,V_2$.
At rank 18 add all energy bounds and Poisson on $\langle9\rangle$.
At rank 24 add the energy congruences, \eqref{eq:ac}, the aggregate
energy congruences, and Poisson on the spaces named
\texttt{V2\_27}, \texttt{V1\_27}, \texttt{V2\_9a},
\texttt{V2\_9b}, \texttt{V3}.
At rank 27 add the seven spaces involving 7 or $V_4$ in Section 4.
Thereafter, each new orbit adds its energy and Poisson on its
representative line. Residual congruences at $0,21$ remain active;
the $63$ congruence is also applied after representatives $3,5,15$.
At the last stage, all 64 energies must equal 4096.
Every Poisson check covers all 64 frequencies and coset characters;
early checks on subsets merely reject failures sooner.

The rank-41 count after orbit 5 represents $23,636\cdot4=94,544$
full coefficient vectors. Expanding these without the remaining
Poisson tests gives 32,176 satisfying energy at $b=0$, 39,520 at
$b=63$, 5,552 at both, and zero at all 64 values of $b$.
This redundancy again uses necessary energy conditions, not an
asserted equivalence with bentness.

\section{Reproduction and trust boundary}
The release contains three Python programs and the execution logs.
The main sieve uses 14 representative fibers and the signed transport
formula. The second sieve uses coefficient vectors in $\F^{43}$,
all 64 fiber truth tables, and a butterfly Walsh transform; it never
uses signed transport to obtain a fiber spectrum. They share the
mathematical filters, so agreement addresses implementation errors
but is not an independent proof that those filters are necessary.
Section 4 supplies that proof.

All arithmetic relevant to acceptance is exact. Walsh values on a
fiber lie in $[-64,64]$, total fiber energy is at most 262,144, and
all such quantities fit the signed 32-bit arrays used by the sieves.
The binary matrix products in unsigned bytes are subsequently reduced
modulo two; reduction modulo 256 preserves their parity. The foundation
checks use signed 64-bit transforms.

The main program tests sixteen quadratic RS bent functions against
the 24 Poisson tables, and checks transport and Poisson identities
against complete Walsh transforms on these and fifty deterministic
random quartic examples. The quadratic controls are not asserted to
pass quartic-specific sector conditions. The foundation script checks
signed transport on all 43 generators, all fibers and frequencies,
as well as complete Walsh exclusion for $n=4,6,8$.

The Lean file proves finite arithmetic identities, the corrected
quarter-period orbit assertion, and the parity syzygy evaluated on
the actual list of 43 monomial orbit representatives. Its compilation
does not formalize completeness of that list, the whole Fourier
argument, the Python execution, or Theorem~\ref{thm:main}.
The release records Lean 4.15.0 and the Python/NumPy versions.

\section{Relation to smaller cases and earlier work}
The case $n=4$ is elementary: the only nonzero homogeneous quartic
function is $x_0x_1x_2x_3$, with $W_f(0)=14$, so it is not bent.
It is a useful convention check, not evidence of a new general theorem.
The supplied broader verifier also reproduces the $n=10$ correction:
the residues of $\pm8$ and $\pm32$ modulo five form the same set.
Its nonzero-frequency derivative checks, not the invalid sign argument,
exclude the remaining branch. These computations are supplementary
and are not required for Theorem~\ref{thm:main}.

Sun et al.~\cite{Sun}, Theorem 2, give necessary orbit conditions under
explicit prime-size hypotheses. Their condition for $n=2p^r$, $p\ge5$,
already covers the odd number of short orbits condition in degree four:
a stabilizer of a four-element support acts freely on that support,
so its order divides both four and $2p^r$. Thus the only short
quartic orbits have length $p^r$, and are bi-antipodal.
For $n=12$ their stated hypothesis with $v_2(n)=2$ would require the
odd prime factor to exceed 16, and therefore does not apply directly.
This comparison with one theorem is not a claim that no other published
result covers twelve variables. The contribution presented here is the
explicit certificate and its auditable spectral reduction.

\paragraph{Data and code.}
Source files, exact output records, a checksum manifest, and reproduction
instructions accompany this manuscript. No external dataset is needed.
\paragraph{Limitations.}
The certificate does not settle any dimension $n\ge14$. Tests on
generators in larger dimensions establish only the specific linear
identities checked; they are not exhaustive bentness searches over
the full coefficient spaces.

\begin{thebibliography}{9}
\bibitem{SM} P. St\u{a}nic\u{a} and S. Maitra,
Rotation symmetric Boolean functions---count and cryptographic applications,
\emph{Discrete Applied Mathematics} 156(10), 1567--1580 (2008).
\bibitem{MCF} Q. Meng, L. Chen and F.-W. Fu,
On homogeneous rotation symmetric bent functions,
\emph{Discrete Applied Mathematics} 158, 1111--1117 (2010).
\bibitem{Sun} L. Sun, Z. Shi, J. Liu and F.-W. Fu,
On the conjecture about the nonexistence of homogeneous rotation symmetric
bent functions, \emph{Designs, Codes and Cryptography} 94, article 93 (2026).
\url{https://doi.org/10.1007/s10623-026-01848-4}.
\bibitem{Carlet} C. Carlet, \emph{Boolean Functions for Cryptography and
Coding Theory}, Cambridge University Press (2021).
\end{thebibliography}
\end{document}
